\textbf{Keywords:} Tapeout, Tiny Tapeout, PTAT, Time-domain readout,
Temperature chamber, Transfer curve, INL, Calibration, Allan deviation,
Quantisation, Dither, Bernoulli statistics, Random telegraph noise

\section{Measured Silicon}\label{measured-silicon}\index{measured silicon@Measured silicon}

The previous chapter set you a project. This one is what happened when
two groups finished it.

In the spring of 2025 two groups of three students designed temperature
sensors in this course, and both went to a shuttle: Tiny Tapeout project
258, \texttt{tt\_um\_jnw\_wulffern}, on ttsky25a in sky130. The chips
came back. This chapter is what they do.

It is here for a specific reason. Everything else in this book is either
theory or somebody else's measurement. This is the one chapter where the
circuit was designed by students at your stage, on the tools you are
using, and then measured against a temperature chamber - and where the
interesting results are not the ones the designers were aiming at.

\subsection{The same physics, read out two
ways}\label{the-same-physics-read-out-two-ways}

Both sensors are the same idea, and it is the idea the project chapter
asks you to build: a current proportional to absolute temperature (PTAT)
charging a capacitor, with a comparator watching the ramp. The current
comes from the difference between two diode voltages at different
current densities, which is the bandgap core of the references chapter,

\[I(T) = \frac{kT}{qR}\ln N\]

and the time to charge \(C\) up to a threshold \(V_{ref}\) is therefore

\[t = \frac{V_{ref}C}{I(T)} \propto \frac{1}{T}\]

The time falls as the temperature rises, so the \emph{rate} \(1/t\)
rises with absolute temperature. Everything measured below is done in
the rate domain for that reason: it is the quantity that should be a
straight line through the origin.

Both groups built that PTAT core the same way - a pnp of unit area
against eight of them, an amplifier forcing the two branch voltages
equal, and a resistor setting the current - but not with the same
circuit. GR07 puts three RPPO16 in series and mirrors the result
straight out, one to one. GR06 uses an RPPO8 and an RPPO4 and then
divides its current down through two further mirrors, an NMOS pair ten
wide and a PMOS pair ten wide, before it reaches the capacitor. That is
roughly a hundredfold division, and it is most of why GR06 runs six
times slower on a capacitor a quarter the size.

Both also take their comparator threshold from a resistive divider
across the supply rather than from anything absolute: GR07 taps one
resistor up from ground in a string of four, so \(V_{ref} = V_{DD}/4\);
GR06 taps one of three, so \(V_{ref} = V_{DD}/3\). Neither threshold is
a bandgap voltage. The ratio \(V_{ref}/I(T)\) that sets the time
therefore carries the supply in it, which is worth remembering before
reading any absolute number below as a property of the sensor alone.

Where the two designs really part company is what closes the loop, and
that turns out to matter far more than anything in the analogue.


{
\centering
\aicfig{media/jnwtt_gr07_tikz.pdf}
}

\emph{Figure 1: GR07, by Reidar Arne Eidsvik Nerheim, Pol Batalle Largo
and Tord Olsen Sætermo, drawn from the taped-out netlist. The comparator
output is buffered into a D flip-flop clocked by the project clock, and
the flip-flop's output both leaves the chip as PWM and turns on the two
NMOS that short the ramp capacitor. The loop closes through the chip, so
the sensor free-runs and its period is the observable}


{
\centering
\aicfig{media/jnwtt_gr06_tikz.pdf}
}

\emph{Figure 2: GR06, by Gabin Sbaffi, Erik K. Jensen and Renate
Klemetsdal. The same front end, and the same NMOS shorting the capacitor
- but its gate is a chip input. The host asserts reset, the capacitor
discharges, and on release the ramp runs once. One pulse per stimulus,
and its width is the observable. Nothing in this circuit is clocked}

At room temperature GR07 free-runs at about 910 kHz - a period near 1.1
µs - so half a second of capture contains half a million periods. GR06
only produces a pulse when the host asks for one, and the host asks
about 4.5 thousand times a second, so the same half second contains
about 2 300 pulses. GR07 therefore delivers roughly two hundred times
more events per measurement, which you would expect to make it the
better sensor. Hold that thought.

\section{What the silicon does}\label{what-the-silicon-does}\index{what the silicon does@What the silicon does}

\subsection{Against a real reference}\label{against-a-real-reference}\index{against a real reference@Against a real reference}

Everything you can do on a bench calibrates a sensor against itself. To
find out whether either sensor is \emph{right}, rather than merely
repeatable, you need an outside opinion. A Vötsch temperature chamber
stepping from 5 °C to 70 °C in 5 K steps, logging its own probe beside
both sensors, supplies one. Each point below is the mean of the last 90
s of a dwell, after both the oven and the die have stopped moving.


{
\centering
\aicfig{media/jnwtt_transfer_tikz.pdf}
}

\emph{Figure 3: Both sensors against the chamber's own probe, each rate
divided by its own value at 25 °C. The dashed line is what a current
strictly proportional to absolute temperature would do. Both lie on it:
the physics works, and the two designs agree with each other despite a
sixfold difference in rate}

Figure 3 is the result the students were designing for, and it is a good
one. Two independently designed PTAT cores, built by different people
from the same principle, land on the ideal line and on each other. GR07
gives 2.98 kHz/K on 907 kHz; GR06 gives 0.48 kHz/K on 142 kHz. Those are
0.33 \%/K and 0.34 \%/K - the same fractional slope, which is what
``proportional to absolute temperature'' means.

Figure 4 is where the two part company. Take the straight line away and
both are left with about 1.3 K of something. Now allow one more term -
the \(T\ln T\) curvature the references chapter warns about, which comes
from the temperature dependence of the saturation current in the very
diodes that make the PTAT current.

For GR06 that removes almost all of it: 1.33 K peak becomes 0.23 K, and
the residual goes flat. GR06's error \emph{is} bandgap curvature, and it
is the textbook shape at the textbook size.

For GR07 the same term barely helps - 1.61 K becomes 1.04 K - and what
is left still wanders. Whatever limits GR07 is not curvature.


{
\centering
\aicfig{media/jnwtt_inl_tikz.pdf}
}

\emph{Figure 4: What is left of each transfer after the best straight
line, and after also allowing a \(T\ln T\) term - the bandgap curvature
the references chapter warns about. Almost all of GR06's residual is
that curvature. Almost none of GR07's is}

One caveat on the identification. Over a 65 K span \(T\ln T\) and
\(T^2\) are nearly collinear, and fitting either removes the same
amount. The data cannot tell you which functional form it is; what it
can tell you is that GR06's residual is smooth curvature of the size and
sign a bandgap predicts, and that GR07's is not smooth curvature at all.

\subsection{What calibration buys}\label{what-calibration-buys}\index{what calibration buys@What calibration buys}

The question a product actually asks is not ``how linear is it'' but
``how many oven visits must I pay for''. Every calibration temperature
costs money in production, so the interesting curve is error against
number of trim points.


{
\centering
\aicfig{media/jnwtt_cal_gr06_tikz.pdf}
}

\emph{Figure 5: GR06's error after calibrating at one, two and three
temperatures. Each extra trim point buys accuracy, ending at ±0.35 K
over 5-70 °C. This is what a well-behaved sensor looks like}


{
\centering
\aicfig{media/jnwtt_cal_gr07_tikz.pdf}
}

\emph{Figure 6: The same for GR07, which gets \emph{worse} going from
one point to two}

GR06 does what you would hope: 1.93 K with one point, 1.37 K with two,
0.35 K with three. That is a normal, well-behaved sensor, and ±0.35 K
over a 65 K span from three trim points is a respectable number for a
first silicon by three students.

GR07 goes 1.96 K, then 2.21 K, then 1.38 K. Getting worse when you give
it more information looks like a mistake in the analysis. It is not. A
two-point calibration corrects a \emph{slope}, and GR07's error is not a
slope. Trimming the line just pivots it and drops the residual somewhere
else.

To see what GR07's error actually is, you have to look at what its
output is made of.

\section{The clock in the loop}\label{the-clock-in-the-loop}\index{the clock in the loop@The clock in the loop}

\subsection{GR07's output is quantised, and noise is what rescues
it}\label{gr07s-output-is-quantised-and-noise-is-what-rescues-it}

Look again at Figure 1. The comparator trips whenever it trips, but its
output reaches the reset transistors through a flip-flop, and a
flip-flop can only change at a clock edge. Every GR07 period is
therefore a whole number of 64 MHz clock cycles, and nothing in between
is representable.

That is not a small quantum. GR07's period is about seventy clock cycles
at room temperature and falls by roughly one cycle for every five kelvin
of warming, so \textbf{one clock cycle is about 4.75 K} - three times
the whole calibrated error of Figure 6, and larger than the raw
uncalibrated error of Figure 4.

A noiseless version of this circuit would be a plain 4.75 K quantiser.
Every period would come out the same, and the average of a million of
them would tell you no more than one of them did: 4.75 K steps, and
nothing in between.

What gets it below that is noise. Jitter on the comparator crossing -
from the PTAT current, from the comparator itself, from the supply - is
enough to push some periods over the clock edge and not others. Say the
crossing sits a fraction \(f\) of a cycle past an edge. The output then
alternates between \(N\) and \(N+1\) cycles, and the fraction of periods
that come out long is an estimate of \(f\). Averaging a million of them
reads it to a few millikelvin.

That is dither: noise deliberately relied upon to make a coarse
quantiser resolve below its own step. It is the same mechanism that lets
a noisy ADC average its way below one LSB, and it is why converters are
sometimes given dither on purpose. Here nobody gave it any - the circuit
simply happened to have enough.

And a mechanism that depends on noise being big enough has an obvious
way to fail.


{
\centering
\aicfig{media/jnwtt_deadzone_tikz.pdf}
}

\emph{Figure 7: GR07's measured noise against \(f\), the fractional part
of its period in 64 MHz clock cycles. The dashed curve is
\(\sqrt{f(1-f)}\) with one scale factor fitted to the data}

The measured noise is not constant: it varies by a factor of eight
across the sweep, and it varies with \(f\) in a very particular way.

That shape is not a coincidence. If a fraction \(f\) of the periods come
out long and the rest short, each period is a coin flip with probability
\(f\) - a Bernoulli trial, whose variance is \(f(1-f)\). The spread of a
rate estimated from many of them therefore goes as \(\sqrt{f(1-f)}\):
smallest where the period is nearly a whole number of cycles, largest
where it sits half-way between. The dashed curve is that expression with
a single scale factor fitted, and it follows the fourteen measurements
with a correlation of 0.997.

So GR07's precision depends on the temperature it happens to be at,
through nothing more than the arithmetic of the quantiser. At 10 °C its
period is within a hundredth of a whole cycle and its noise is 22 Hz; at
25 °C it sits nearly half-way between and the noise is 186 Hz. Same
circuit, same clock, eight times the noise, and it is the fractional
part of a division that decides which you get.

This also puts a floor under how well GR07 can ever be trimmed.
Calibration fits a smooth function to a transfer; what GR07 delivers is
a 4.75 K staircase whose dither statistics change from step to step. The
cheapest real fix is not more calibration points but a faster project
clock, which shrinks the step in proportion.

GR06, with no clock anywhere in its path, has none of this. Its output
is a pulse width in continuous time, and it is limited by something else
entirely.

\subsection{When averaging stops
helping}\label{when-averaging-stops-helping}

If the resolution comes from averaging, the honest question is how long
you can usefully average for. A standard deviation cannot answer it. The
Allan deviation can: it asks how repeatable the answer is if you average
for \(\tau\) seconds, and its \emph{slope} is the information. Falling
means averaging still buys you precision. Rising means drift has taken
over and averaging longer makes the answer worse.


{
\centering
\aicfig{media/jnwtt_allan_tikz.pdf}
}

\emph{Figure 8: Allan deviation of both sensors over fifteen minutes in
a quiet room. GR06 is lower at every averaging time. GR07 bottoms out at
57 mK after about eight seconds and then degrades; GR06 keeps improving
out past a minute, to 31 mK}

The two-hundred-fold event advantage does not appear anywhere on this
plot. GR06 is the better sensor at every averaging time, from the
shortest to the longest, and the gap widens.

That is worth sitting with, because within a single half-second capture
GR07 really is the more precise of the two: its standard error is 3.3 mK
against GR06's 12.8 mK. But note how little that is. Two hundred times
the events should be \(\sqrt{200}\), a factor of fifteen, and GR07 only
manages four. The missing factor is its own single-event resolution: one
GR07 period is worth 2.3 K of scatter against 0.6 K for one GR06 pulse,
because 4.75 K of quantiser has to be dithered through. GR07 spends most
of its event advantage buying back what the flip-flop took.

The four it keeps then does not survive either. Its Allan deviation at
the shortest averaging time is 80 mK - twenty-four times its own
statistical error - where GR06's is 55 mK against 12.8 mK, a factor of
four. Both sensors are limited by something other than counting
statistics, and GR07 by far the more so. Beyond about ten seconds its
curve turns upward, which means something slower than the averaging
window is now spoiling the answer and averaging longer makes it worse.
GR06, producing two hundred times fewer events, keeps improving.

There is a second result hiding in the same run. Over those fifteen
minutes the two sensors' fluctuations are essentially uncorrelated
(\(r = 0.11\)) despite sitting on the same die, in the same room,
sharing a supply. If the wander were the room, both would see it - so it
looks like each sensor's own noise rather than the room.

A quiet room cannot prove that on its own, though: ``nothing happened''
and ``both sensors missed it'' look identical from inside. The next
section is the experiment that separates them.

\subsection{Two sensors, one event}\label{two-sensors-one-event}\index{two sensors, one event@Two sensors, one event}

Everything so far measures a sensor against a reference. The two sensors
on this die can also be measured against each other, which asks a
different question: not \emph{what is the temperature} but \emph{did
something happen}. For that you need something to happen - so, a can of
freeze spray, and then a finger held on the package.


{
\centering
\aicfig{media/jnwtt_spray_tikz.pdf}
}

\emph{Figure 9: A can of freeze spray at 13 s, then a fingertip on the
package from 88 s to 122 s, both sensors on one capture. The die falls
13 K in under two seconds, at better than 20 K/s at the steepest. Gaps
are dead time between captures, drawn as gaps rather than interpolated}

The two traces do not sit on top of each other - through the long
recovery GR06 reads up to about a kelvin above GR07 - but they move
together. Every feature appears in both at the same instant, and the
correlation of the two temperature series is 0.96 through the spray and
recovery and 0.99 while the finger is on.

The offset is not a surprise. Each sensor is calibrated at a single
point near 23 °C against the ideal line through the origin, with no
offset term to absorb comparator delay or reset time, so at 8 °C both
are extrapolating fifteen kelvin beyond their only anchor. What the
figure supports is the \emph{shape}: the timing, the rates, and the fact
that both saw the same thing. The absolute depth is the weakest number
on the page.

This is the control the previous section needed. Quiet, the two sensors
were uncorrelated; given something real to follow, they agree almost
perfectly. Both cannot be true of the same physical temperature, and
that settles it: the wander in the quiet room was not the room, it was
each sensor's own noise. A real thermal event moves both. Noise moves
one.

Which is the argument for putting two of anything on a die. One sensor
gives you a number and no way to know whether to believe it. Two give
you a way to tell a measurement from an artefact.

Now the same experiment, gently: four breaths on the package.


{
\centering
\aicfig{media/jnwtt_breath_tikz.pdf}
}

\emph{Figure 10: Four breaths, both sensors, one capture. Each breath
moves both - but GR06 reads about 2.2 times the excursion GR07 does}

Both sensors agree that four things happened, and disagree about how big
they were by a factor of 2.2. The obvious reading is that the two
circuits have different gains. They do not: over the 15 K of Figure 9
they agree to within a few per cent, and a real gain error would show up
\emph{more} strongly over a wider excursion, not vanish.

So something compresses GR07's response to small, fast signals without
touching its response to large slow ones. A one-kelvin breath is about a
fifth of GR07's 4.75 K quantiser step, so the quantiser is the obvious
suspect - but the chamber sweep of Figure 7 shows its dither behaving
the way ideal dither should, which would give a linear response at any
excursion. The two observations are not yet reconciled, and nothing else
measured here settles it.

What can be said without hand-waving is the practical part: for small,
fast signals the two sensors disagree by a factor of about two, and GR06
is the one that agrees with the larger excursion. Trust it. Establishing
\emph{why} would need a controlled small-step experiment with GR07
parked at several known values of \(f\) - which is a good project, and
is not in this data set.

\section{One trap}\label{one-trap}\index{one trap@One trap}

\subsection{GR06 does not drift, it
switches}\label{gr06-does-not-drift-it-switches}

GR06 is the better-behaved sensor everywhere above, so it is worth
asking what actually limits \emph{it}. Hold the chamber still and look
at a minute of GR06 with the slow dwell drift removed.


{
\centering
\aicfig{media/jnwtt_rts_tikz.pdf}
}

\emph{Figure 11: Sixty seconds of GR06 at a fixed chamber temperature,
with the slow dwell drift removed. The reading sits flat, drops about
six tenths of a kelvin to a second level, stays there a second or two,
and comes back}

Thermal noise would give a fuzzy band around zero. This gives two
levels. That is random telegraph noise - a single charge trap in the
silicon capturing and emitting one carrier - and it is the same
mechanism the noise chapter introduces as the microscopic origin of
flicker noise, here big enough to see one trap at a time.

The step is about 0.6 K from 5 °C to 55 °C - it drifts from 0.7 K to 0.4
K across that span, but it does not scale with anything - which says it
is one trap rather than a population. The time it stays trapped is
thermally activated and follows an Arrhenius law, which is what
identifies it as a trap rather than as something in the instrument.


{
\centering
\aicfig{media/jnwtt_rts_life_tikz.pdf}
}

\emph{Figure 12: The trap's mean lifetime in the low state against
inverse thermal energy, one point per chamber dwell. A straight line
here is an Arrhenius law; the slope is an activation energy of 226 meV.
Above about 55 °C the two levels merge into the noise and the fit stops
meaning anything, so those dwells are left out}

So GR06's real single-reading accuracy is not its ±0.35 K three-point
calibration - that is a ninety-second mean. A single reading can be six
tenths of a kelvin off whenever the trap happens to be occupied. Unlike
GR07's quantiser this is cheap to fix: a median over a second or two
steps straight over it.

\section{Summary}\label{summary}

Two sensors from this course, on the same die, from the same principle.

\begin{itemize}
\tightlist
\item
  Both PTAT cores work. The rates follow absolute temperature to within
  the residual of Figure 3, and the two designs - different resistors,
  and a hundredfold current division in one of them - agree with each
  other to a per cent in fractional slope.
\item
  GR06, with three calibration points, is accurate to ±0.35 K over 5-70
  °C. What is left after a straight line is bandgap curvature: a
  \(T\ln T\) term takes 1.33 K down to 0.23 K. Its limit for a single
  reading is not linearity at all but one charge trap worth about 0.6 K,
  which a short median removes.
\item
  GR07 has four times GR06's statistical precision within a capture and
  is the worse sensor at every averaging time - its readings scatter
  twenty-four times more than its own statistics predict. Its output is
  quantised to whole clock cycles, 4.75 K each, and only noise dithering
  the crossing lets averaging resolve below that. How precise it is then
  depends on where it sits in that step: the noise follows
  \(\sqrt{f(1-f)}\) to a correlation of 0.997, an eightfold range across
  the sweep, decided by the fractional part of a division.
\item
  Calibration cannot absorb a staircase. A faster clock shrinks it.
\item
  The two sensors' noise is uncorrelated on one die. Two sensors tell
  you when something is real; one tells you a number.
\end{itemize}

If you take one thing from this chapter into your own project: both
groups' analogue cores did what they were designed to do. Every awkward
result came from the boundary - a flip-flop, a clock frequency, a choice
of what to observe. That is where to spend your attention.

\section{Would you like to know
more?}\label{would-you-like-to-know-more}

The measurement setup, the full analysis and every figure's underlying
data, in far more detail than fits here \citeproc{ref-jnwtt25}{{[}1{]}}

The Tiny Tapeout shuttle that carried it, and how to put your own design
on one \citeproc{ref-tinytapeout}{{[}2{]}}

Where the PTAT core comes from, and why its curvature is not what limits
these parts \citeproc{ref-razavi21}{{[}3{]}}

Quantisation, dither and why a coarse converter with noise on its input
beats a coarse converter without \citeproc{ref-schreier.uds}{{[}4{]}}

Random telegraph noise as the microscopic origin of flicker noise
\citeproc{ref-kirton89}{{[}5{]}}

\phantomsection\label{refs}
\begin{CSLReferences}{0}{0}
\bibitem[\citeproctext]{ref-jnwtt25}
\CSLLeftMargin{{[}1{]} }%
\CSLRightInline{C. Wulff, {``{JNW-TEMP}: Two temperature sensors,
measured.''} 2026 {[}Online{]}. Available:
\url{https://analogicus.com/jnw-tt-2025/presentation.html}}

\bibitem[\citeproctext]{ref-tinytapeout}
\CSLLeftMargin{{[}2{]} }%
\CSLRightInline{M. Venn and the Tiny Tapeout contributors, {``Tiny
tapeout.''} {[}Online{]}. Available: \url{https://tinytapeout.com}}

\bibitem[\citeproctext]{ref-razavi21}
\CSLLeftMargin{{[}3{]} }%
\CSLRightInline{B. Razavi, {``The design of a low-voltage bandgap
reference {[}the analog mind{]},''} \emph{IEEE Solid-State Circuits
Magazine}, vol. 13, no. 3, pp. 6--16, 2021, doi:
\href{https://doi.org/10.1109/MSSC.2021.3088963}{10.1109/MSSC.2021.3088963}.
}

\bibitem[\citeproctext]{ref-schreier.uds}
\CSLLeftMargin{{[}4{]} }%
\CSLRightInline{R. Schreier and G. C. Temes, \emph{Understanding
delta-sigma data converters}. John Wiley \& Sons, 2005. }

\bibitem[\citeproctext]{ref-kirton89}
\CSLLeftMargin{{[}5{]} }%
\CSLRightInline{M. J. Kirton and M. J. Uren, {``Noise in solid-state
microstructures: A new perspective on individual defects, interface
states and low-frequency ({1/f}) noise,''} \emph{Advances in Physics},
vol. 38, no. 4, pp. 367--468, 1989, doi:
\href{https://doi.org/10.1080/00018738900101122}{10.1080/00018738900101122}.
}

\end{CSLReferences}
